\section{Discussion}
\label{sec:discussion}

\subsection{What the numbers mean for a foundry}

Width and thickness monitors made of strips cannot report the sidewall angle of an oxide-clad full etch, however
many they are and in whichever polarization. The cost of that blindness is a misprediction of rib devices of
\qty{710.8}{pm} per degree of tilt. That is comparable to the \qty{841.5}{pm} within-die spread that Zortman \textit{et~al.}\
reported for microdisks, whose across-wafer spread was \qty{8.014}{nm}~\cite{zortman2010}, and well above the
\qty{\pm 32}{pm} accuracy reached by trimming~\cite{jayatilleka2021}. A bound of
\qty{0.39}{\degree} on the angle sounds coarse next to these numbers. What a designer uses, however, is the predicted
resonance of other devices, and within the shared-parameter model the two-width set predicts a strip ring to
\qty{0.457}{pm} and a rib ring to \qty{0.773}{pm} at a \qty{1}{pm} readout (Section~\ref{sec:prediction}). The angle
is uncertain, but the combination of angle, width and thickness that sets each resonance is not. That prediction
needs an absolute wavelength reference to about \qty{10}{pm}; without one it degrades to \qty{982.2}{pm}.

\subsection{When the shared-parameter assumption breaks}

The central limitation is the assumption, shared by every multi-device extraction~\cite{nrc2026}, that all
structures of a set see the same process. Section~\ref{sec:variation} shows that the angle extraction is far more
sensitive to it than a width or thickness extraction. The bound of the two-width set doubles once the structures'
widths differ by \qty{0.03938}{nm} or their angles by \qty{0.01749}{\degree}, and uncorrected offsets of
\qty{0.25}{\degree} produce biases of tens of degrees. For comparison, the best linewidth control reported for a
CMOS line is \qty{0.6}{nm} within a chip~\cite{selvaraja2010}, and the mismatch between identical rings grows with
their separation~\cite{chrostowski2014}, with layout-dependent systematic components~\cite{xing2023}. The closest
measured number is from Xing \textit{et~al.}, who extracted nominally identical waveguides at 44 sites on one die. Their width
had a standard deviation of \qty{1.26}{nm}, and after removing a linear trend across the die the local variation
between nearby devices still reached \qty{3}{nm}, against \qty{0.5}{nm} for the thickness, with an extraction error of
\qty{0.37}{nm} in width~\cite{xing2018}. These
are differences between copies of one design, and the quantity that matters here, the difference between adjacent
waveguides of different widths after the shared offset is removed, has not been reported. If it is of the size Xing
\textit{et~al.}\ measured, the shared-parameter bound is not reachable. A per-structure width spread of
\qty{0.5}{nm} already raises the two-width bound to \qty{3.38}{\degree} (Fig.~\ref{fig:variation}(b)).

A working angle test structure therefore has to meet a consistency requirement that width monitors do not face. Its
members must be close enough, and similar enough in local pattern density, that their departures from one another
stay below a few hundredths of a nanometre, or the design must measure those departures. Two consequences follow for
design. First, the two-width set tolerates about ten times larger departures than the best set of eight, so the
smaller set is the better proposal even though its shared-parameter bound is higher. Second, linewidth depends on local pattern density~\cite{xing2023}, and if the angle does too it can
change systematically with waveguide width; a trend of \qty{0.25}{\degree} per
\qty{100}{nm} biases the two-width estimate by \qty{6.869}{\degree}. Width ranges should therefore be kept narrow, or
the trend added to the model as one more parameter.

\subsection{Relation to physical metrology}

Cross-sections and AFM measure the profile directly at a few sites and remain the reference for calibrating any
optical extraction~\cite{orji2018,yoo2022}. Scatterometry is fast but model-based~\cite{raymond2005}, and its targets
are designed with the same Fisher-information tools we use here~\cite{klein2025,silver2007}. The optical structures
studied here measure the effect of the angle on the guided modes, which is the quantity a circuit depends on, at
every place where rings can be placed. They need a calibrated cross-section model to turn that effect into degrees,
and the co-location requirement above to make the conversion meaningful.

\subsection{Limitations}

The results are simulations, and no device was fabricated or measured. The symmetry argument is first-order and holds
for oxide above and below the core. An air-clad strip is not mirror-symmetric (its angle/width ratio in our solves is
0.9782), and we computed no bound for air-clad sets. The walls are ideal trapezoids tilting about their bottom edge,
without corner rounding or roughness, the material indices are nominal values without dispersion, and $\ngr$ uses a fixed
\qty{5}{nm} wavelength step; none of these choices was varied. Rings are assumed to have known resonance orders and
lengths, so that each resonance gives $\neff$ directly. Only a common calibration offset was tested, and an unknown
length per ring was not. Our $\neff$ noise at \qty{1}{pm} is below the fitting error of measured MZIs
(Section~\ref{sec:theory}), so the measurement model may be optimistic for $\neff$. The bounds scale linearly with
$\sigma_\lambda$ and depend on the noise conversion by tens of
percent. Halving or doubling the rib noise relative to the strips changes the two-width bound by
\qty{-47.9}{\percent} and \qty{97.5}{\percent}, and the ring length used in the conversion changes it by
\qty{-49.4}{\percent} at \qty{50}{\micro m} and \qty{97.8}{\percent} at \qty{200}{\micro m}. The priors matter much
less, at most \qty{0.6}{\percent} for factor-of-two changes. The two-width bound carries a numerical uncertainty of
about \qty{5}{\percent}, the change when the domain is enlarged. The best set of eight is limited by derivative accuracy,
moving by \qty{-15.3}{\percent} between remeshed and fixed-mesh derivatives and by \qty{-27.5}{\percent} at half the
derivative step, and fits of femwell data spread 1.487 times its bound (Section~\ref{sec:surrcheck}). We treat the rib and strip angles as independent because the two etches are separate steps, and
Section~\ref{sec:results} shows how much a link would help. The analysis covers TE ribs only and the
\qty{1.55}{\micro m} band.

\subsection{Next steps}

The decisive experiment is a test chip that carries the two-width set in a compact layout next to conventional width
monitors, measured at controlled temperature against an absolute wavelength reference and compared with
cross-sections of the same dies. Repeating the set at several separations on the chip would measure the local
departures that Section~\ref{sec:variation} identifies as the limit. On the modeling side, per-structure departures
could be absorbed by structures whose members share most of their waveguide, or by the single-ring, many-resonance
approach of~\cite{nrc2026} extended to a rib. TM ribs, air cladding and other platforms such as silicon nitride, where
the cladding and etch profile differ, remain open.
